% Figure 3 from the drawing board: three approaches to "the spin factors follow the labels, up to the statistics sign; parity is separate".
\documentclass[10pt,border=2mm]{standalone}
\usepackage[T1]{fontenc}
\usepackage{figurestyle}
\input{primitives.tex}
\input{molecules.tex}
\input{numbers.tex}
\newcommand\spinup[2]{\draw[line width=.7pt,draw=PlateInk,-{Stealth[length=1.4mm,width=1.1mm]}] (#1,#2-.17)--(#1,#2+.17);}
\newcommand\spindn[2]{\draw[line width=.7pt,draw=PlateInk,-{Stealth[length=1.4mm,width=1.1mm]}] (#1,#2+.17)--(#1,#2-.17);}
\begin{document}
\begin{tikzpicture}[plate]
\path[use as bounding box] (0,0) rectangle (15,20.5);
% ================================================= A. spin arrows on the nuclei; the digits move, the arrows stay; the sign on the arrow
\node[lab,anchor=west] at (.2,20.1) {A. spin arrows beside the nuclei: relabel moves the digits, the arrows stay where they are; the exchange sign rides on the arrow};
\begin{scope}[shift={(2.2,17.6)}]\molsolid\mollabelsinside\pic[scale=1.0] at (0,0) {mol3d-water-X};
  \spinup{-1.28}{.58}\spindn{1.0}{.37}\node[lab] at (0,-1.2) {$X$};\end{scope}
\draw[harrow] (3.9,17.6)--(5.3,17.6); \node[lab,anchor=south] at (4.6,17.68) {$\sigma=(12)$}; \node[sub,anchor=north] at (4.6,17.5) {$\times\,\stat(\sigma)=-1$};
\begin{scope}[shift={(7.0,17.6)}]\molsolid\mollabelsinside\def\molA{2}\def\molB{1}\pic[scale=1.0] at (0,0) {mol3d-water-X};
  \spinup{-1.28}{.58}\spindn{1.0}{.37}\node[lab] at (0,-1.2) {$\sigma X$};\end{scope}
\node[sub,anchor=north west,text width=5.4cm] at (9.3,18.6) {$\Psi(\sigma X)\sim\stat(\sigma)\,\sigma\act\Psi(X)$: the same points of space carry the same spins; only the digits moved, and a proton exchange costs the sign};
% ================================================= B. kets beside the nuclei, the product value written on the molecule
\node[lab,anchor=west] at (.2,15.3) {B. the value written on the molecule as kets beside the nuclei};
\begin{scope}[shift={(2.2,12.9)}]\molsolid\mollabelsinside\pic[scale=1.0] at (0,0) {mol3d-water-X};
  \node[lab,anchor=east] at (-1.2,.95) {$|{\uparrow}\rangle$}; \node[lab,anchor=west] at (.95,.75) {$|{\downarrow}\rangle$}; \node[lab,anchor=north] at (0,-.5) {$|\omega\rangle$}; \node[lab] at (0,-1.2) {$X$};\end{scope}
\draw[harrow] (3.9,12.9)--(5.3,12.9); \node[lab,anchor=south] at (4.6,12.98) {$\sigma=(12)$};
\begin{scope}[shift={(7.0,12.9)}]\molsolid\mollabelsinside\def\molA{2}\def\molB{1}\pic[scale=1.0] at (0,0) {mol3d-water-X};
  \node[lab,anchor=east] at (-1.2,.95) {$|{\uparrow}\rangle$}; \node[lab,anchor=west] at (.95,.75) {$|{\downarrow}\rangle$}; \node[lab,anchor=north] at (0,-.5) {$|\omega\rangle$}; \node[lab] at (0,-1.2) {$\sigma X$};\end{scope}
\node[lab,anchor=west] at (9.3,13.3) {$\Psi(X)=|{\uparrow}\rangle_1|{\downarrow}\rangle_2|\omega\rangle_3$};
\node[lab,anchor=west] at (9.3,12.6) {$\Psi(\sigma X)\sim-\,|{\downarrow}\rangle_1|{\uparrow}\rangle_2|\omega\rangle_3$};
% ================================================= C. parity: the point reflection drawn, X and -X in one frame
\node[lab,anchor=west] at (.2,10.6) {C. parity as a drawing: $X$ and $-X$ in one frame, each nucleus joined to its image through the origin; no triad};
\begin{scope}[shift={(4.0,7.6)}]
  \fill[PlateInk] (0,0) circle[radius=1.3pt]; \node[sub,anchor=north west,inner sep=1pt] at (.05,-.05) {$0$};
  \foreach \x/\y in {-0.9348/0.5826,0.6528/0.3684,0.0178/-0.0599} {\draw[line width=.4pt,draw=PlateInk,dash pattern=on .8pt off 1.4pt] ({1.3*\x},{1.3*\y})--({-1.3*\x},{-1.3*\y});}
  \begin{scope}\molsolid\mollabelsinside\pic[scale=1.3] at (0,0) {mol3d-water-X};\end{scope}
  \begin{scope}\molsolid\mollabelsinside\pic[scale=1.3] at (0,0) {mol3d-water-minusX};\end{scope}
  \node[lab] at (-1.5,1.3) {$X$}; \node[lab] at (1.5,-1.3) {$-X=E^*X$};
\end{scope}
\node[sub,anchor=north west,text width=5.6cm] at (8.0,8.6) {$E^*$ sends every column through the origin; the digits and the spins come along unchanged, and the value picks up the parity $\pm1$ instead of $\stat$. Water is planar, so $-X$ is also a rotated copy};
% ================================================= D. the triad on a sphere, redrawn with the two-tone shading and octant shading after Harter 5.1.1
\node[lab,anchor=west] at (.2,4.9) {D. the handedness glyph redrawn: a tetrahedral corner (three edges from one vertex) and its inversion; readable without a sphere};
\begin{scope}[shift={(3.0,2.4)}]
  \draw[heavy] (0,0)--(0,1.4); \draw[heavy] (0,0)--(1.3,-.35); \draw[heavy] (0,0)--(-.7,-.75);
  \node[sub,anchor=south] at (0,1.45) {$z$}; \node[sub,anchor=west] at (1.35,-.35) {$y$}; \node[sub,anchor=north east] at (-.7,-.78) {$x$};
  \fill[PlateInk] (0,0) circle[radius=1.5pt];
  \node[sub,anchor=north,text width=3cm,align=center] at (.2,-1.1) {right-handed: $x\times y=z$};
\end{scope}
\draw[harrow] (5.4,2.4)--(6.6,2.4); \node[lab,anchor=south] at (6.0,2.48) {$E^*$};
\begin{scope}[shift={(9.0,2.4)}]
  \draw[heavy] (0,0)--(0,-1.4); \draw[heavy] (0,0)--(-1.3,.35); \draw[heavy] (0,0)--(.7,.75);
  \node[sub,anchor=north] at (0,-1.45) {$-z$}; \node[sub,anchor=east] at (-1.35,.35) {$-y$}; \node[sub,anchor=south west] at (.7,.78) {$-x$};
  \fill[PlateInk] (0,0) circle[radius=1.5pt];
  \node[sub,anchor=north,text width=3.4cm,align=center] at (-.2,-1.6) {left-handed: $(-x)\times(-y)=z\neq-z$; no rotation does this};
\end{scope}
\end{tikzpicture}
\end{document}
